\begin{table}[H]
  \centering
  \caption{Configuración del sistema de envasado}
  \label{tab:anexoA-configuracion}

  \begin{tabular*}{\textwidth}{@{\extracolsep{\fill}}>{\raggedright\arraybackslash}p{0.32\textwidth} c c >{\raggedright\arraybackslash}p{0.40\textwidth}}
    \toprule
    \textbf{Parámetro} & \textbf{Valor} & \textbf{Unidad} & \textbf{Comentario} \\
    \midrule
    Líneas de producción & 1 & línea & Realiza embotellado y etiquetado, acoplados o desacoplados \\
    Capacidad de la línea por período & 84 & horas & Tiempo disponible para producción y set-ups \\
    Períodos del horizonte & 3 & períodos & Coinciden con las etapas del árbol de escenarios \\
    Estanques intermedios disponibles & 20 & estanques & Número máximo utilizable en cada período \\
    Capacidad de cada estanque & 10.000 & litros & Un estanque alimenta la línea de embotellado \\
    Tamaño de botella (todos los vinos) & 0,75 & litros & $\approx$ 13.333 botellas por estanque \\
    Subutilización máxima de un estanque & 100\% & fracción & No se exige uso mínimo del estanque \\
    \bottomrule
  \end{tabular*}
\end{table}

\begin{table}[H]
  \centering
  \caption{Productos, etiquetas admisibles e inventarios iniciales}
  \label{tab:anexoA-productos}

  % Panel A
  \begin{tabular*}{\textwidth}{@{\extracolsep{\fill}}>{\raggedright\arraybackslash}p{0.20\textwidth} >{\raggedright\arraybackslash}p{0.35\textwidth} >{\raggedright\arraybackslash}p{0.38\textwidth}}
    \multicolumn{3}{@{}l}{\textbf{Panel A.} \textit{Productos y etiquetas admisibles}} \\[0.4em]
    \toprule
    \textbf{Vino} & \textbf{Etiquetas admisibles} & \textbf{Productos terminados (vino, etiqueta)} \\
    \midrule
    Vino 1 & Etiquetas 1, 2 y 3 & (1,1), (1,2), (1,3) \\
    Vino 2 & Etiquetas 1 y 2    & (2,1), (2,2) \\
    \bottomrule
  \end{tabular*}

  \vspace{1.2em}

  % Panel B
  \begin{tabular*}{\textwidth}{@{\extracolsep{\fill}}>{\raggedright\arraybackslash}p{0.55\textwidth} c >{\raggedright\arraybackslash}p{0.30\textwidth}}
    \multicolumn{3}{@{}l}{\textbf{Panel B.} \textit{Inventario inicial}} \\[0.4em]
    \toprule
    \textbf{Inventario inicial} & \textbf{Valor} & \textbf{Unidad} \\
    \midrule
    Botellas llenas sin etiquetar (ambos vinos)    & 0 & botellas \\
    Productos terminados (todas las combinaciones) & 0 & botellas \\
    Pedidos atrasados (\textit{backlog}) iniciales   & 0 & botellas \\
    \bottomrule
  \end{tabular*}

  \vspace{0.6em}
  \raggedright \small \textit{Nota.} La demanda del período 1 (nodo 1) se satisface con actividades de embotellado y etiquetado predefinidas, iguales a la demanda de ese nodo.
\end{table}

\begin{table}[H]
  \centering
  \caption{Tiempos y costos de operación}
  \label{tab:anexoA-tiempos-costos}

  % Panel A
  \begin{tabular*}{\textwidth}{@{\extracolsep{\fill}}>{\raggedright\arraybackslash}p{0.52\textwidth} c r}
    \multicolumn{3}{@{}l}{\textbf{Panel A.} \textit{Tiempos y costos de preparación (set-ups)}} \\[0.4em]
    \toprule
    \textbf{Concepto} & \textbf{Tiempo (h)} & \textbf{Costo (\$)} \\
    \midrule
    Set-up de embotellado & 1,5 & 1.500 \\
    Set-up conjunto de embotellado y etiquetado & 1,5 & 1.500 \\
    Set-up de etiquetado & 0,5 & 500 \\
    \bottomrule
  \end{tabular*}

  \vspace{1.2em}

  % Panel B
  \begin{tabular*}{\textwidth}{@{\extracolsep{\fill}}>{\raggedright\arraybackslash}p{0.52\textwidth} c r}
    \multicolumn{3}{@{}l}{\textbf{Panel B.} \textit{Tiempos unitarios de proceso}} \\[0.4em]
    \toprule
    \textbf{Proceso} & \textbf{Valor (h/botella)} & \textbf{Unidad} \\
    \midrule
    Llenado (solo embotellar) & 0,00014 & h/botella \\
    Etiquetado (solo etiquetar) & 0,00028 & h/botella \\
    Llenado y etiquetado acoplados & 0,00028 & h/botella \\
    \bottomrule
  \end{tabular*}

  \vspace{1.2em}

  % Panel C
  \begin{tabular*}{\textwidth}{@{\extracolsep{\fill}}>{\raggedright\arraybackslash}p{0.52\textwidth} c r}
    \multicolumn{3}{@{}l}{\textbf{Panel C.} \textit{Costos de inventario y servicio}} \\[0.4em]
    \toprule
    \textbf{Costo de inventario y servicio} & \textbf{Valor} & \textbf{Unidad} \\
    \midrule
    Mantener una botella llena sin etiquetar & 10 & \$/botella$\cdot$período \\
    Mantener una botella de producto terminado & 15 & \$/botella$\cdot$período \\
    Pedido atrasado (\textit{backorder}) & 15.000 & \$/botella$\cdot$período \\
    \bottomrule
  \end{tabular*}
\end{table}

\begin{table}[H]
  \centering
  \caption{Árbol de escenarios de demanda}
  \label{tab:anexoA-arbol-escenarios}

  \begin{tabular*}{\textwidth}{@{\extracolsep{\fill}}cccccc}
    \toprule
    \textbf{Nodo} & \textbf{Período} & \textbf{Estado} & \textbf{Nodo antecesor} & \textbf{Prob. condicional} & \textbf{Prob. del nodo} \\
    \midrule
    1  & 1 & Media & — & 1,000 & 1,000 \\
    2  & 2 & Alta  & 1 & 0,300 & 0,300 \\
    3  & 2 & Media & 1 & 0,400 & 0,400 \\
    4  & 2 & Baja  & 1 & 0,300 & 0,300 \\
    5  & 3 & Alta  & 2 & 0,700 & 0,210 \\
    6  & 3 & Media & 2 & 0,300 & 0,090 \\
    7  & 3 & Alta  & 3 & 0,333 & 0,133 \\
    8  & 3 & Media & 3 & 0,333 & 0,133 \\
    9  & 3 & Baja  & 3 & 0,333 & 0,133 \\
    10 & 3 & Media & 4 & 0,600 & 0,180 \\
    11 & 3 & Baja  & 4 & 0,400 & 0,120 \\
    \bottomrule
  \end{tabular*}

  \vspace{0.6em}
  \raggedright \small \textit{Nota.} Árbol de 3 períodos y 11 nodos. La probabilidad de cada nodo es el producto de las probabilidades condicionales desde la raíz; las probabilidades de los nodos terminales suman 1.
\end{table}

\begin{table}[H]
  \centering
  \caption{Demanda por producto y nodo}
  \label{tab:anexoA-demanda}

  \small
  \setlength{\tabcolsep}{3pt} % Reduce el espacio horizontal entre columnas
  \begin{tabular*}{\textwidth}{@{\extracolsep{\fill}} l *{11}{r} }
    \toprule
    \textbf{Producto} & \textbf{N1} & \textbf{N2} & \textbf{N3} & \textbf{N4} & \textbf{N5} & \textbf{N6} & \textbf{N7} & \textbf{N8} & \textbf{N9} & \textbf{N10} & \textbf{N11} \\
    \midrule
    (1,1) & 23.525 & 53.218 & 34.328 & 18.118 & 81.856 & 45.302 & 50.426 & 31.482 & 12.538 & 30.856 & 10.426 \\
    (1,2) & 24.328 & 15.184 & 28.525 & 45.832 & 24.546 & 31.371 & 23.569 & 40.767 & 57.965 & 20.546 & 45.569 \\
    (1,3) & 20.368 & 38.150 & 21.080 & 6.762  & 4.920  & 25.765 & 35.043 & 25.646 & 16.248 & 10.920 & 11.043 \\
    (2,1) & 21.080 & 39.387 & 30.368 & 12.646 & 34.115 & 22.773 & 37.045 & 27.831 & 18.616 & 24.115 & 17.045 \\
    (2,2) & 22.991 & 38.048 & 27.991 & 14.949 & 33.440 & 22.824 & 35.085 & 25.147 & 15.209 & 23.440 & 15.085 \\
    \bottomrule
  \end{tabular*}

  \vspace{0.6em}
  \raggedright \small \textit{Nota.} Demanda expresada en botellas. La demanda de un nodo debe atenderse en ese nodo; lo no atendido queda como pedido atrasado. Las demandas de los productos del Vino 1 presentan correlaciones negativas entre sí; las del Vino 2, positivas (ver Sección 4.4).
\end{table}

\begin{table}[H]
  \centering
  \caption{Matriz comparativa de metodologías de IO para el proceso de envasado y postergación}
  \label{tab:comparativa_metodologias}

  \small
  \setlength{\tabcolsep}{3pt} % Controla el espacio entre columnas para evitar desbordes
  \begin{tabular*}{\textwidth}{@{\extracolsep{\fill}}
    >{\raggedright\arraybackslash}p{0.15\textwidth} 
    >{\raggedright\arraybackslash}p{0.18\textwidth} 
    >{\raggedright\arraybackslash}p{0.18\textwidth} 
    >{\raggedright\arraybackslash}p{0.18\textwidth} 
    >{\raggedright\arraybackslash}p{0.16\textwidth}}
    \toprule
    \textbf{Método} & \textbf{Manejo de Incertidumbre} & \textbf{Manejo de \textit{Set-ups} y Tiempos} & \textbf{Ventaja Principal} & \textbf{Desventaja Principal} \\
    \midrule

    \textbf{Optimización Estocástica (2 Etapas)} & Árboles discretos / probabilidades condicionales \cite{gupta2000two}. & \textit{Set-ups} de embotellado en Etapa 1; WIP y penalizaciones en Etapa 2. & Cuantifica el valor del \textit{postponement}, VSS y \textit{risk pooling} ($r = 0{,}66$). & Explosión combinatoria de escenarios si aumenta el árbol. \\
    \addlinespace[0.5em]

    \textbf{MILP Determinista (Valor Esperado)} & Ignora aleatoriedad; usa demandas promediadas \cite{gunther2014block}. & Excelente para \textit{set-ups} dependientes de la secuencia y filtros \cite{lillo2025optimal}. & Alta velocidad de cómputo en \textit{solvers} como Gurobi. & Inflexible; ignora variabilidad y subestima pedidos atrasados. \\
    \addlinespace[0.5em]

    \textbf{Programación Robusta (ARO)} & Conjuntos poliédricos; optimiza para el peor escenario \cite{hu2020hybrid}. & Introduce \textit{time gaps} entre corridas de producción \cite{lappas2016multi}. & Inmune a variaciones extremas; no requiere de distribuciones previas. & ``Precio de la Robustez'' elevado; inventario WIP conservador e innecesario. \\
    \addlinespace[0.5em]

    \textbf{Simulación Monte Carlo / DES} & Muestreo estocástico de trayectorias \cite{review2024simula}. & Representación detallada de fallas y tiempos no lineales \cite{jeon2016survey}. & Excelente para pruebas de estrés y validación operacional. & Herramienta puramente descriptiva, no prescriptiva; no genera planes óptimos. \\
    \bottomrule
  \end{tabular*}
\end{table}

\begin{table}[H]
  \centering
  \caption{Dimensiones consideradas en los escenarios experimentales}
  \label{tab:escenarios_metodologicos}

  \small
  \setlength{\tabcolsep}{4pt}
  \begin{tabular*}{\textwidth}{@{\extracolsep{\fill}}
    >{\raggedright\arraybackslash}p{0.25\textwidth} 
    >{\raggedright\arraybackslash}p{0.30\textwidth} 
    >{\raggedright\arraybackslash}p{0.40\textwidth}}
    \toprule
    \textbf{Dimensión} & \textbf{Niveles} & \textbf{Interpretación} \\
    \midrule
    Variabilidad de demanda
    & $10\%$, $30\%$, $50\%$
    & Intensidad de la incertidumbre respecto de la demanda media. \\
    \addlinespace[0.4em]

    Capacidad de línea
    & 21, 42, 63 y 84 h/período
    & Grado de holgura disponible para absorber producción, recurso y \textit{set-ups}. \\
    \addlinespace[0.4em]

    Nivel de postergación
    & $PL \in \{0,2,3,4,5\}$
    & Número máximo de etiquetas para las cuales se permite mantener flexibilidad mediante diferenciación tardía. \\
    \addlinespace[0.4em]

    Correlación de demanda
    & Positiva / negativa
    & Grado en que la demanda de productos se mueve conjuntamente o presenta compensaciones entre etiquetas. \\
    \bottomrule
  \end{tabular*}
\end{table}